\clearpage
\VerdierPhysicalPageAnchor{0269}%
\begingroup
\small
\setlength{\parindent}{0pt}%
\setlength{\parskip}{0pt}%
\setlength{\emergencystretch}{3em}%
\newcommand{\VerdierContentsRow}[3]{%
  \par\noindent
  \hangindent=#1\relax\hangafter=1
  \hspace*{#1}#2\nobreak
  \leaders\hbox{\kern.32em.\kern.32em}\hfill
  \nobreak\hbox{\kern.65em #3}\par}
\newcommand{\VerdierContentsChapter}[1]{%
  \par\addvspace{.55\baselineskip}\noindent\textbf{#1}\par\smallskip}
\newcommand{\VerdierContentsSection}[2]{%
  \par\addvspace{.45\baselineskip}\noindent
  \hfill\textbf{#1.}\enspace #2\hfill\mbox{}\par\smallskip}

\VerdierSemanticAnchor{toc-ch2}%
\VerdierContentsChapter{Chapitre II : Catégories triangulées.}
\VerdierSemanticAnchor{toc-ch2-sec1}%
\VerdierContentsSection{1}{Définitions et premières propriétés.}
\VerdierSemanticAnchor{toc-ch2-sec1-1}%
\VerdierContentsRow{.5em}{\textbf{1.1.}\enspace Définition des catégories triangulées}{93}
\VerdierSemanticAnchor{toc-ch2-sec1-2}%
\VerdierContentsRow{.5em}{\textbf{1.2.}\enspace Premières propriétés des catégories triangulées}{97}
\VerdierSemanticAnchor{toc-ch2-sec1-3}%
\VerdierContentsRow{.5em}{\textbf{1.3.}\enspace Exemples de catégories triangulées}{106}

\VerdierSemanticAnchor{toc-ch2-sec2}%
\VerdierContentsSection{2}{Localisation dans les catégories triangulées.}
\VerdierSemanticAnchor{toc-ch2-sec2-1}%
\VerdierContentsRow{.5em}{\textbf{2.1.}\enspace Systèmes multiplicatifs de morphismes}{111}
\VerdierSemanticAnchor{toc-ch2-sec2-2}%
\VerdierContentsRow{.5em}{\textbf{2.2.}\enspace Construction de la catégorie localisée}{117}
\VerdierSemanticAnchor{toc-ch2-sec2-3}%
\VerdierContentsRow{.5em}{\textbf{2.3.}\enspace Propriétés du foncteur de localisation}{125}

\VerdierSemanticAnchor{toc-ch2-sec3}%
\VerdierContentsSection{3}{Abélianisation des catégories triangulées.}
\VerdierSemanticAnchor{toc-ch2-sec3-1}%
\VerdierContentsRow{.5em}{\textbf{3.1.}\enspace Construction de la catégorie $\mathrm A(\mathcal D)$}{135}
\VerdierSemanticAnchor{toc-ch2-sec3-2}%
\VerdierContentsRow{.5em}{\textbf{3.2.}\enspace Propriétés de la catégorie $\mathrm A(\mathcal D)$}{137}

\VerdierSemanticAnchor{toc-ch2-sec4}%
\VerdierContentsSection{4}{Objets spectraux.}
\VerdierSemanticAnchor{toc-ch2-sec4-1}%
\VerdierContentsRow{.5em}{\textbf{4.1.}\enspace Définition des objets spectraux}{145}
\VerdierSemanticAnchor{toc-ch2-sec4-2}%
\VerdierContentsRow{.5em}{\textbf{4.2.}\enspace Mécanisme des suites spectrales}{147}
\VerdierSemanticAnchor{toc-ch2-sec4-3}%
\VerdierContentsRow{.5em}{\textbf{4.3.}\enspace Les termes $E\!\begin{bmatrix}n\\-\infty,p,q,+\infty\end{bmatrix}$}{153}
\VerdierSemanticAnchor{toc-ch2-sec4-4}%
\VerdierContentsRow{.5em}{\textbf{4.4.}\enspace Convergence}{155}

\VerdierSemanticAnchor{toc-ch3}%
\VerdierContentsChapter{Chapitre III : Les catégories dérivées.}
\VerdierSemanticAnchor{toc-ch3-sec1}%
\VerdierContentsSection{1}{Définitions des catégories dérivées.}
\VerdierSemanticAnchor{toc-ch3-sec1-1}%
\VerdierContentsRow{.5em}{\textbf{1.1.}\enspace Notations et terminologie}{157}
\VerdierSemanticAnchor{toc-ch3-sec1-2}%
\VerdierContentsRow{.5em}{\textbf{1.2.}\enspace Les sous-catégories remarquables des catégories dérivées}{159}
\VerdierSemanticAnchor{toc-ch3-sec1-3}%
\VerdierContentsRow{.5em}{\textbf{1.3.}\enspace Le $\delta$-foncteur canonique}{164}

\VerdierSemanticAnchor{toc-ch3-sec2}%
\VerdierContentsSection{2}{Résolutions.}
\VerdierSemanticAnchor{toc-ch3-sec2-1}%
\VerdierContentsRow{.5em}{\textbf{2.1.}\enspace Morphismes élémentaires}{169}
\VerdierSemanticAnchor{toc-ch3-sec2-2}%
\VerdierContentsRow{.5em}{\textbf{2.2.}\enspace Étude de certains ensembles d'objets}{173}
\VerdierSemanticAnchor{toc-ch3-sec2-3}%
\VerdierContentsRow{.5em}{\textbf{2.3.}\enspace Quelques lemmes sur les résolutions}{177}
\VerdierSemanticAnchor{toc-ch3-sec2-4}%
\VerdierContentsRow{.5em}{\textbf{2.4.}\enspace Applications}{185}

\VerdierSemanticAnchor{toc-ch3-sec3}%
\VerdierContentsSection{3}{Étude des Ext.}
\VerdierSemanticAnchor{toc-ch3-sec3-1}%
\VerdierContentsRow{.5em}{\textbf{3.1.}\enspace Différentes définitions équivalentes}{191}
\endgroup
